\textbf{1. Kiến thức cần nhớ}
\begin{itemize}
\item \textbf{Biến cố}: Là một sự kiện liên quan đến phép thử, được biểu diễn bởi một tập con $A \subset \Omega$.
\item \textbf{Kết quả thuận lợi}: Các phần tử thuộc tập hợp $A$ gọi là các kết quả thuận lợi cho biến cố $A$. Số kết quả thuận lợi kí hiệu là $n(A)$.
\item \textbf{Công thức tính xác suất cổ điển} (khi các kết quả có cùng khả năng xảy ra):
$$\boxed{P(A) = \dfrac{n(A)}{n(\Omega)} = \dfrac{\text{Số kết quả thuận lợi cho } A}{\text{Số phần tử của không gian mẫu}}}$$
\item \textbf{Tính chất}:
\begin{itemize}
\item $0 \le P(A) \le 1$.
\item Xác suất của biến cố chắc chắn bằng $1$: $P(\Omega) = 1$.
\item Xác suất của biến cố không thể bằng $0$: $P(\emptyset) = 0$.
\item Biến cố đối $\overline{A}$: $\boxed{P(\overline{A}) = 1 - P(A)}$.
\end{itemize}
\end{itemize}

\textbf{2. Phương pháp giải các dạng}

%% DANG 9D8?26-1
\textbf{Dạng 1. Xác định các kết quả thuận lợi cho biến cố}
\begin{enumerate}
\item Đọc kĩ điều kiện của biến cố $A$ (ví dụ: tổng số chấm bằng $7$, rút được quả bóng chẵn/nguyên tố).
\item Lọc từ không gian mẫu $\Omega$ ra các phần tử thỏa mãn điều kiện để lập tập hợp $A$.
\item Đếm số phần tử $n(A)$.
\end{enumerate}

%% DANG 9D8?26-2
\textbf{Dạng 2. Tính xác suất của biến cố}
\begin{enumerate}
\item Tính số phần tử không gian mẫu $n(\Omega)$.
\item Tính số kết quả thuận lợi $n(A)$.
\item Lập tỉ số $P(A) = \dfrac{n(A)}{n(\Omega)}$ và rút gọn thành phân số tối giản hoặc số thập phân.
\end{enumerate}

%% DANG 9D8?26-3
\textbf{Dạng 3. Bài toán xác suất thực tế}
\begin{enumerate}
\item Nhận diện mô hình thực tế: Rút thăm trúng thưởng, bốc thăm ngẫu nhiên học sinh, kiểm tra chất lượng sản phẩm, chia bảng thi đấu thể thao\dots
\item Xác định tổng số đối tượng tham gia ($n(\Omega)$) và số đối tượng đạt tiêu chuẩn ($n(A)$).
\item Tính xác suất và trả lời bài toán.
\end{enumerate}
